\section{Limitations}
\label{sec:limitations}

\subsection{The map is not compact, and the size is load-bearing}
\label{sec:lim:compact}
Our maps carry $25$--$100\times$ more primitives than the GS-SLAM baselines at
$8$--$16\times$ the peak GPU memory. The natural defence would be that this is
redundancy --- prune it and parity remains. We tested that defence and it is
false.

\Cref{tab:curve} truncates our own refined maps to successively smaller budgets
by keeping the highest-opacity Gaussians (the ones the optimiser is itself
fading are the honest ones to drop first), re-scoring on the identical
held-out frames with no re-optimisation afterwards.

\begin{table}[t]
\centering\small
\setlength{\tabcolsep}{4pt}
\begin{tabular}{@{}lrrrr@{}}
\toprule
Gaussians kept & office0 & office1 & room0 & room2 \\
\midrule
full ($\sim$2--3M)     & 26.29 & 22.08 & 25.44 & 23.62 \\
1{,}000{,}000          & 24.68 & 21.71 & 25.01 & 22.87 \\
300{,}000              & 20.23 & 18.96 & 16.26 & 17.87 \\
83{,}372               & 13.54 & 13.97 & \phantom{0}8.68 & 10.43 \\
23{,}019               & 12.04 & 12.08 & \phantom{0}5.16 & \phantom{0}7.74 \\
\midrule
\emph{Photo-SLAM at $\sim$83K} & \emph{22.23} & \emph{17.80} & \emph{17.92} & \emph{20.75} \\
\bottomrule
\end{tabular}
\caption{PSNR versus map budget, four Replica scenes. At Photo-SLAM's own
primitive budget we are $8.7$/$3.8$/$9.2$/$10.3$\dB{} \emph{worse} than it,
having been $4.08$/$4.28$/$7.53$/$2.87$\dB{} better at $25$--$36\times$ the
count. The knee sits consistently near $1$M Gaussians.}
\label{tab:curve}
\end{table}

We state the bound in both directions. This is naive post-hoc truncation
without re-optimisation, so every row is a \emph{lower} bound --- a system
actually trained or refined at that budget would do meaningfully better. But
the caveat does not rescue the defence: the measurement shows our quality is
\emph{constituted by} having millions of primitives rather than padded by them.
Training or refining directly at a small budget, or a densification-control
policy, is future work we have not done.

\subsection{Scope of the comparison}
Rendering on TUM is $n{=}1$ and cannot be extended, because MonoGS's monocular
tracking diverges on the remaining sequences. Optimisation budgets are matched
by wall clock on TUM but not on Replica, where Photo-SLAM runs to its own
shutdown iteration. Resolution and undistortion handling still differ per
system and we have not quantified their residual effect. Photo-SLAM was built
against OpenCV~4.10 rather than the 4.7/4.8 its authors tested, because our
CUDA~12.4 requirement forces it; this touches its image front-end but not its
Gaussian mapping.

\subsection{Unresolved threads}
Three questions we could not close, stated rather than omitted. The
cross-family variation in head-only gains survived four falsified attributions
and remains unexplained. Cross-frame colour consistency is diagnosed but not
fixed --- our photospatial augmentation happens to reduce hue disagreement by
$12.6\%$ on the sequence that originally exposed the bug, which is a
side-effect rather than a solution, since the single-view objective contains
nothing that penalises cross-view inconsistency. Refitting the loop-closure
retrieval assets on Splatt3R features looked favourable offline (Recall@1
$+20\%$ relative) but regressed catastrophically on one of three sequences in
real SLAM, and was rejected --- one more instance of an offline metric win that
did not survive the online check.
